% C1 Inicio del Preámbulo
% C1.1 Definición de la Clase de Documento
\documentclass[12pt, doc]{apa7}

% C1.2 Configuración de Idioma y Fuentes
\usepackage[spanish]{babel}
\usepackage{fontspec}
\setmainfont{Times New Roman}

% C1.3 Geometría y Maquetación de la Página
\usepackage[left=1in,right=1in,top=1in,bottom=1in]{geometry}
\usepackage{setspace}
\parindent=0.5in

% C1.4 Elementos Visuales y Manejo de Color
\usepackage{graphicx} % <--- Soporte para la inserción de imágenes externas (PNG, JPG, PDF)
\usepackage{xcolor} % <--- Soporte para la gestión y creación de colores personalizados
\definecolor{codegray}{rgb}{0.95,0.95,0.95} % <--- Color gris claro para el fondo de los bloques de código
\definecolor{codeblue}{rgb}{0.0,0.4,0.7} % <--- Color azul oscuro corporativo para palabras clave

% C1.5 Hipervínculos y Navegación del Documento
\usepackage{csquotes}
\usepackage{hyperref}
\hypersetup{
	pdfborder={0 0 0},
	colorlinks=true,
	linkcolor=black,
	urlcolor=blue,
	citecolor=black
}

% C1.6 Estructura de Tablas Avanzadas
\usepackage{array}
\usepackage{tabularx}
\usepackage{ltablex}
\keepXColumns
\usepackage{booktabs}

% C1.7 Bloques de Código Fuente (Listings)
\usepackage{listings} % <--- Paquete principal para incrustar y dar estilo a fragmentos de código
\lstset{
	backgroundcolor=\color{codegray}, % <--- Aplica el fondo gris personalizado
	basicstyle=\ttfamily\small, % <--- Tipografía monoespaciada y tamaño ligeramente más pequeño
	breaklines=true, % <--- Ajusta automáticamente las líneas que superan los márgenes del texto
	keywordstyle=\color{codeblue}, % <--- Aplica el color azul a las palabras clave nativas
	frame=single, % <--- Dibuja un recuadro simple de una sola línea alrededor del código
	rulecolor=\color{lightgray}, % <--- Aplica un color gris claro al borde del recuadro
	numbers=left, % <--- Activa la numeración de líneas en el margen izquierdo
	numberstyle=\tiny\color{gray}, % <--- Estilo gris y diminuto para los números de línea
	stepnumber=1, % <--- Incrementa los números de línea de uno en uno (numera todas las líneas)
	showstringspaces=false % <--- Desactiva símbolos extraños en los espacios en blanco de cadenas de texto
}

% C1.8 Definición de Lenguaje YAML Personalizado para Listings
\lstdefinelanguage{yaml}{
	keywords={version, services, db, image, container_name, restart, environment, ports, volumes}, % <--- Palabras clave de Docker
	keywordstyle=\color{codeblue}\bfseries, % <--- Resalta las palabras clave en azul y negrita
	basicstyle=\ttfamily\small, % <--- Fuerza el uso de la fuente monoespaciada
	sensitive=false, % <--- Ignora la distinción estricta entre mayúsculas y minúsculas
	comment=[l]{\#}, % <--- Define el carácter '#' como inicio de comentario de una línea
	commentstyle=\color{gray}\ttfamily, % <--- Aplica color gris a los comentarios
	stringstyle=\color{purple}\ttfamily, % <--- Aplica color púrpura al texto entre comillas
	moredelim=[l][\color{orange}]{\:} % <--- Resalta los dos puntos (:) de YAML en color naranja
}



% C1.9 Metadatos y Afiliación Académica
\title{}
\author{}
\shorttitle{}
\affiliation{
	SENA - Centro de Electricidad, Electrónica y Telecomunicaciones \\
	Tgo. en Gestión e Implementación de Bases de Datos \\
	\today
}

% Configuración de bibliografía
\usepackage[style=apa, backend=biber]{biblatex}
\addbibresource{referencias.bib}

% ---------------------- FIN DEL PREAMBULO --------------------------------------------------